% TOP_PROBLEM: {"id": 3323, "title": "Crossing numbers and coloring", "category_id": 3, "category": {"id": 3, "name": "graph_theory", "display_name": "Graph Theory", "description": "Problems involving graphs, networks, and their properties.", "slug": "graph-theory", "order_index": 3, "created_at": "2026-07-31T15:26:25.671Z"}, "status": "open", "problem_number": "OPG-37068", "set_id": 12}
% FIRST_POSTED: 2026-10-02
% ATTEMPT_STATUS: unresolved
% MODEL: GPT 6 Astra Ultra
% UnsolvedMath ID 3323; source catalogue code OPG-37068
% !TeX program = lualatex
% Research attempt dated 2026-10-02; canonical statement preserved.
\documentclass[11pt,a4paper]{article}
\usepackage[margin=25mm]{geometry}
\usepackage{fontspec}
\setmainfont{lmroman10-regular.otf}[BoldFont=lmroman10-bold.otf,
 ItalicFont=lmroman10-italic.otf,BoldItalicFont=lmroman10-bolditalic.otf]
\usepackage{amsmath,amssymb,mathtools}
\usepackage{unicode-math}
\setmathfont{latinmodern-math.otf}
\AtBeginDocument{\let\setminus\smallsetminus}
\usepackage{xurl,hyperref}
\hypersetup{colorlinks=true,urlcolor=blue}
\setcounter{secnumdepth}{0}
\setlength{\parindent}{0pt}
\setlength{\parskip}{5pt}
\setlength{\emergencystretch}{3em}
\begin{document}
\section*{Crossing numbers and coloring}\label{3323}
{\small UnsolvedMath ID 3323 · OPG-37068 · Graph Theory · OpenGarden}
\subsection{Definitions and mathematical statement}
Let $G$ be a finite simple undirected graph. Its chromatic number
$\chi(G)$ is the least number of colours in a vertex colouring
where adjacent vertices have different colours. A plane drawing
represents vertices by distinct points and edges by simple arcs
between their ends, with no edge through another vertex, finitely
many transverse crossings of edge interiors, and no triple crossing.
The crossing number $\operatorname{cr}(G)$ is the least total number
of crossing points over all such drawings; endpoints are not counted.
Edges need not be straight.

For an integer $t\ge1$, write $K_t$ for the complete graph on $t$
vertices. Albertson's conjecture is the universal implication
\[
                          \chi(G)\ge t
                 \quad\Longrightarrow\quad
                          \operatorname{cr}(G)\ge\operatorname{cr}(K_t).
\]
The right-hand side is the actual crossing number of $K_t$, not a
conjectured formula for that number. Thus defining this question
does not require solving the separate complete-graph crossing problem.

Equivalently, among graphs with chromatic number exactly $t$, the
complete graph should minimize crossing number. To pass to the
``at least'' formulation, delete vertices until a subgraph of
chromatic number $t$ remains, and use monotonicity of crossing number
under taking subgraphs. The conjecture does not assert that $G$
contains a complete subgraph, a complete subdivision, or a complete
minor. It compares two numerical invariants and permits arbitrarily
many vertices and arbitrary graph structure beyond the chromatic
hypothesis.
\subsection{Short English statement}
Does every graph requiring at least t colours have crossing number at least that of the complete graph on t vertices?
\subsection{Research attempt}
\paragraph{Scope and method.}
This research attempt by GPT-6 Astra Ultra is dated 2 October 2026.
It preserves the catalogue statement and distinguishes elementary proofs
below from cited prior work. No novelty claim or formal proof-assistant
verification is made. The final paragraph states the precise remaining scope.
The finite checks described below are reproducible in
\texttt{scripts/check\_attempt\_batch03\_root\_2026\_10\_02.py}.
\paragraph{A route through critical graphs.}
The target compares a chromatic obstruction with the crossing
number of a clique, rather than with an arbitrary polynomial in
the chromatic number. We derive a sufficient condition based on
the order of a critical subgraph. It proves the desired inequality
for graphs whose critical obstruction is sufficiently large,
but leaves a concrete range of smaller obstructions untreated.
This is an elementary reduction, not a claim to improve the
literature surveyed in [S2].

\paragraph{An elementary crossing estimate.}
Let $H$ be simple with $n$ vertices, $m$ edges and crossing number
$c$. Deleting at most one edge for each crossing in an optimal
drawing leaves a planar graph. The planar edge bound, including
the cases with fewer than three vertices, therefore gives the
convenient weak inequality $c\ge m-3n$.
We may choose an optimal drawing without crossings between
incident edges: exchanging the initial portions of two such
edges removes their crossing without increasing the total.
Keep each vertex independently with probability $p$, and use
the induced drawing. Each edge survives with probability $p^2$,
and each original crossing with probability $p^4$, since its
four endpoints are distinct. Applying the weak inequality to
each outcome and taking expectations gives
\[
 p^4c\ge p^2m-3pn.
\]
If $m\ge4n>0$, choose $p=4n/m\le1$ and rearrange to obtain
\[
                         c\ge\frac{m^3}{64n^2}.
\]
The drawing inherited by a vertex subset need not be optimal;
its crossing count still exceeds the minimum and hence obeys
the same lower estimate. This avoids an unjustified exchange
of minimization and expectation.

\paragraph{A proved range of the conjectured comparison.}
Suppose $\chi(G)\ge t$. Delete vertices until an induced
subgraph $H$ is minimal with chromatic number at least $t$.
Removing one vertex decreases chromatic number by at most one,
so $\chi(H)=t$ and every proper vertex deletion is
$(t-1)$-colourable. A vertex of degree at most $t-2$ could be
added back using an available colour; therefore
$\delta(H)\ge t-1$ and $m\ge(t-1)n/2$.
For $t\ge9$ the preceding estimate applies and yields
\[
 \operatorname{cr}(G)\ge\operatorname{cr}(H)
 \ge\frac{(t-1)^3n}{512}.
\]
Place $t$ points in convex position and draw the complete graph
with straight segments. Each four-point subset contributes
exactly one crossing pair. A generic perturbation avoids triple
crossings, giving $\operatorname{cr}(K_t)\le\binom t4$.
Consequently the required comparison is proved whenever
\[
       n\ge\frac{512\binom t4}{(t-1)^3},\qquad t\ge9.
\]
The condition concerns the critical subgraph's order. Adding
isolated vertices to $G$ cannot make it true, because they do
not enlarge that obstruction.

\paragraph{What remains after this reduction.}
Any counterexample with $t\ge9$ would have a $t$-critical
subgraph with order strictly below this threshold. The threshold
is asymptotic to $(64/3)t$, so the reduction is a linear bound
on the order of a possible critical counterexample for each
fixed chromatic number. It is not a uniform finite verification
of the conjecture as $t$ varies.
For $1\le t\le4$ the desired inequality is immediate because
$K_t$ is planar. The present density calculation by itself
does not settle $5\le t\le8$, nor the remaining small critical
graphs at larger $t$. An attempt to remove the order condition
would need substantially more structural information than
minimum degree: at $n=t$ the displayed lower coefficient is
far below the convex-drawing upper coefficient.

\paragraph{Verification and final status.}
The root script checks the threshold comparisons with exact
rational arithmetic for $9\le t\le100$, including the first
integer order that satisfies each threshold. The probability
and criticality arguments supply the parameter proofs; these
arithmetic checks are not a search over critical graphs.
The sufficient-order theorem is established, but Albertson's
full inequality remains unresolved by this attempt.

\subsection{Sources}
\begin{itemize}
\item[S1] ULAM.AI and the UnsolvedMath Contributors, supplied problems.json, record 3323 (OPG-37068), OpenGarden collection. Catalogue identity and source transcription; CC BY 4.0.
\url{https://huggingface.co/datasets/ulamai/UnsolvedMath}.
\item[S2] Open Problem Garden, Crossing numbers and coloring. Original problem and discussion; the mathematical formulation below is expanded from the supplied source record.
\url{http://www.openproblemgarden.org/op/crossing_numbers_and_coloring}.

\end{itemize}
\end{document}
